\chapter{Introducción}

\section{Antecedentes}
La depresión afecta a más de 280 millones de personas en el mundo y es una de las
principales causas de discapacidad según la OMS \cite{who_depression_2025}. Eso la
pone en el centro de la investigación farmacológica: si no se pueden probar nuevos
medicamentos de forma confiable, el desarrollo de tratamientos se frena.

Antes de pasar a ensayos en humanos, los fármacos se prueban en modelos animales. Uno
de los más usados en roedores es la prueba de nado forzado (Forced Swim Test, FST),
propuesta por Porsolt et al.\ en 1977 \cite{porsolt1977}. En términos simples: se mete
a la rata en un cilindro con agua del que no puede salir y se registra cuánto tiempo
pasa quieta versus cuánto tiempo intenta escapar. Los antidepresivos reducen el tiempo
que el animal pasa inmóvil \cite{slattery2012,armario2021}.

El problema es que analizar esos videos sigue siendo manual: alguien tiene que ver la
grabación y cronometrar cada conducta para cada animal. Eso toma horas y los resultados
dependen del criterio de quien lo hace, así que dos evaluadores distintos pueden llegar
a números diferentes \cite{armario2021}.

Existen herramientas como EthoVision XT \cite{noldus_ethovision} o ezTrack
\cite{pennington2019_eztrack} que automatizan parte del proceso, pero tienen
limitaciones concretas: algunas son demasiado caras, otras solo detectan si el animal
se mueve sin distinguir qué tipo de movimiento hace, y la mayoría requiere condiciones
de grabación más controladas que las del laboratorio real.

El Laboratorio de Neurociencia Conductual de la ENMyH-IPN, dirigido por el Dr. Israel
Salas Ramírez, usa el FST regularmente y tiene un acervo de videos que hoy se analizan
a mano. De esa necesidad concreta nace este Trabajo Terminal.

\section{Planteamiento del problema}
El laboratorio colaborador corre experimentos de nado forzado de forma periódica:
grupos de ratas se evalúan para medir el efecto de distintas moléculas. Cada
experimento genera videos donde se graban cuatro animales al mismo tiempo desde arriba,
con teléfono o tableta, en dos sesiones: 20 minutos el primer día y 5 minutos el
segundo.

Hoy ese análisis se hace completamente a mano: un investigador reproduce el video y va
anotando con un cronómetro cuándo la rata nada, cuándo está inmóvil y cuándo intenta
escalar. El problema tiene varias caras:

\begin{itemize}
  \item \textbf{Tiempo:} analizar un experimento completo a mano puede llevar varias
        horas, según el número de animales y sesiones.
  \item \textbf{Subjetividad:} dos personas pueden usar criterios distintos para decidir
        cuándo empieza o termina una conducta, lo que afecta la reproducibilidad.
  \item \textbf{Sin registro centralizado:} los resultados se guardan en archivos
        separados (hojas de cálculo, notas), sin forma fácil de comparar sesiones o
        grupos entre sí.
  \item \textbf{Análisis temporal limitado:} un desglose minuto a minuto sería útil
        para ver la dinámica del efecto farmacológico, pero hacerlo a mano es
        prácticamente inviable.
\end{itemize}

Todo esto limita cuántos experimentos puede procesar el laboratorio y dificulta
estandarizar sus resultados.

\section{Propuesta de solución}
La propuesta es desarrollar un prototipo de sistema web que use inteligencia artificial
para analizar automáticamente los videos del FST que genera el laboratorio.

El sistema tiene dos partes:

\begin{itemize}
  \item \textbf{Pipeline de análisis de video:} recibe el video, identifica a cada rata
        dentro de su cilindro y clasifica cuadro a cuadro qué conducta realiza (nado,
        inmovilidad o escalamiento) con un modelo de aprendizaje supervisado. La salida
        es el tiempo total de cada conducta por animal y, de forma opcional, el desglose
        por minuto.
  \item \textbf{Plataforma web:} interfaz desde el navegador donde el investigador puede
        subir los videos, ver el estado del análisis, consultar los resultados por animal
        y por sesión, y descargar reportes. Los datos se almacenan en una base de datos
        para facilitar la comparación entre experimentos.
\end{itemize}

El sistema asume que los videos tienen un formato y calidad mínimos: vista superior con
los cuatro cilindros visibles, contraste suficiente entre el animal y el fondo, y
cámara estable. Los videos disponibles hasta ahora tienen variaciones de iluminación y
contraste que pueden afectar el desempeño; durante el desarrollo se evaluará qué tan
robustos son los algoritmos ante esas condiciones y se documentarán los requisitos
mínimos de calidad.

Para validar el sistema, sus resultados se compararán con los registros conductuales
manuales proporcionados por el laboratorio (gold standard), usando métricas de
concordancia que se definirán en la etapa de análisis.

\section{Objetivo general}
Diseñar, desarrollar y validar un prototipo de sistema web asistido por inteligencia
artificial que automatice la detección y cuantificación de las conductas de la prueba
de nado forzado en ratas ---nado activo, inmovilidad e intentos de escape--- con
resultados comparables a la evaluación manual del laboratorio.

\subsection{Objetivos específicos}
\begin{enumerate}
  \item Revisar trabajos y herramientas existentes para el análisis automatizado del
        FST, e identificar los enfoques de visión por computadora y clasificación
        supervisada más adecuados al problema.
  \item Describir el protocolo experimental del FST que usa el laboratorio, incluyendo
        las conductas de interés y las características de los videos, a partir de
        entrevistas técnicas con el equipo del laboratorio y de los registros
        conductuales que éste proporcionará.
  \item Diseñar la arquitectura del sistema, definiendo los módulos de carga de videos,
        pipeline de análisis, base de datos y visualización de resultados.
  \item Implementar el módulo de análisis conductual con visión por computadora y un
        clasificador supervisado para las tres conductas del FST.
  \item Validar el desempeño del sistema comparando sus predicciones contra los
        registros conductuales manuales proporcionados por el laboratorio (gold
        standard), con métricas de concordancia previamente definidas.
  \item Medir la reducción de tiempo que aporta el sistema frente al análisis manual, y
        documentar bajo qué condiciones podría adaptarse a otros experimentos
        conductuales.
\end{enumerate}

\section{Justificación}
El análisis manual del FST tiene limitaciones claras: toma mucho tiempo, depende del
criterio del evaluador y no hay forma centralizada de guardar y comparar los resultados.
Un sistema automatizado resolvería esos tres problemas para el laboratorio colaborador.

Las herramientas comerciales (como EthoVision) son caras y no están pensadas para el
flujo de trabajo del laboratorio. Las de código abierto disponibles no clasifican las
tres conductas del FST ni tienen plataforma web. Este sistema cubre las dos cosas al
mismo tiempo.

Desde el punto de vista académico, el proyecto permite aplicar visión por computadora,
aprendizaje automático y desarrollo web en un problema real con impacto en investigación
biomédica. Eso encaja directamente con el perfil de egreso de la ESCOM.

\section{Alcance del proyecto}

\subsection*{Incluido en el alcance}
\begin{itemize}
  \item Análisis de videos del FST del laboratorio de la ENMyH-IPN, grabados desde
        arriba con cuatro cilindros simultáneos, bajo condiciones de grabación
        suficientes para que el procesamiento sea viable.
  \item Clasificación de tres conductas: nado activo, inmovilidad e intentos de escape.
        Si el sistema no logra distinguir con precisión suficiente nado y escalamiento,
        ambos se reportarán juntos como movilidad.
  \item Métricas por animal y por sesión: tiempo total de cada conducta en los 5 minutos
        del día 2, y desglose por minuto como funcionalidad opcional.
  \item Comparación entre sesiones día 1 y día 2 para el mismo animal o grupo.
  \item Plataforma web con carga de videos, consulta de resultados y generación de
        reportes.
  \item Validación contra el gold standard del laboratorio con métricas de concordancia
        (Kappa de Cohen) y error absoluto medio.
  \item Documentación técnica conforme a los lineamientos de la CATT-ESCOM.
\end{itemize}

\subsection*{Fuera del alcance}
\begin{itemize}
  \item Análisis en tiempo real sobre video en vivo.
  \item Corrección automática de videos con mala iluminación, reflejos excesivos u otros
        problemas de calidad.
  \item Cualquier intervención en los experimentos con animales; eso lo hace el
        laboratorio.
  \item Despliegue institucional definitivo; el producto entregado es un prototipo
        funcional.
  \item Análisis de paradigmas conductuales distintos al FST en esta versión del sistema.
\end{itemize}

\section{Metodología}
El proyecto usa una metodología ágil basada en Scrum, adaptada a un equipo de dos
personas con entregas en dos etapas académicas (TT1 y TT2). La razón de elegir este
enfoque es que el proyecto tiene partes muy distintas entre sí ---el módulo de video,
el clasificador, la web, la base de datos--- y no tiene sentido diseñar todo desde el
inicio si no sabemos cómo va a quedar el módulo de seguimiento hasta probarlo.

Los sprints duran entre 2 y 4 semanas. Al final de cada uno hay un entregable concreto
que se puede revisar con el director.

% TODO: insertar tabla de sprints

Las herramientas son: Python con Flask para el backend y el pipeline; OpenCV y
TensorFlow/scikit-learn para procesamiento de video y clasificación; PostgreSQL para
la base de datos; React.js para el frontend; BORIS para anotar el gold standard; y
Git con Docker para control de versiones y despliegue.